% modB_koopman.tex — module B: the command-conditioned Koopman dynamics model.
% v3 (2026-09-19, post CP2-redo): eq:koopman_model is stated directly on the
% reading window s_t (no encoder/decoder in the main text), so the main chain
% now runs left to right entirely in s-space:
%   s_t -> [K, (L+1)x(L+1)] -> (+) -> shat_{t+1} -> [C selector] -> yhat_{t+1}
% with the command column B (+ the scalar c_t) entering the sum from below --
% this branch is still the point of the figure -- and the bias b entering
% from above. The one thing kept from the lifting story (3.1 P6) is drawn as
% a single faded, optional detour arcing over the top of the main chain,
% through a de-emphasised encoder/decoder pair; the paper reports the
% identity lifting (the solid chain below), so the arc+nets are opacity 0.45
% and never the thing a reader's eye lands on first. The two labels that
% belong to the arc (E_phi, and the "optional learned lifting" caption) are
% kept at full opacity so they stay legible (G4) -- only the drawn shapes
% (line + nets) are faded, per the plan's "opacity=0.45" on the arc itself.
% Canvas 6.53 x 3.03; top-left strip is the title bar (added in right_panel).
\path (0,0) rectangle (6.53,3.03);

% ---------------- 1. finite-memory state s_t --------------------------
% top-left y=2.09, same local coordinate module A uses for its own s_t
% column bottom -- the cross-module arrow in right_panel.tex depends on
% this and is NOT changed here. L+1 is already labelled once, in A; per
% the plan's "no dimension appears twice" rule this module does not repeat
% \DimBraceR for it.
\MemColV{(0.70,2.09)}{4}{0.22}
\node[mathlab,anchor=north] at (0.81,1.19) {$\symMem$};
\draw[thinflow] (0.92,1.75) -- (1.45,1.75);

% ---------------- 2. the transition matrix K ---------------------------
\MatBlock{(1.45,2.19)}{4}{4}{0.22}{clLatentD}
\node[mathlab,anchor=north] at (1.89,1.29) {$\symK$};
\DimBraceT{(1.45,2.21)}{(2.33,2.21)}{$\symMemDim\times{}\symMemDim$}
% the "{}" above is a no-op in the typeset output (same as v20's render);
% it exists only so the G2 audit script's naive string-substitution macro
% expander does not read "\times" immediately followed by \symMemDim's
% own expansion "L+1" as the single (nonexistent) control word "\timesL".
% Real TeX tokenizes \times and \symMemDim as separate control words
% regardless, so this changes nothing LaTeX actually renders -- it is the
% same class of audit-script false positive the plan's risk table already
% flags for \symResid's \lVert, fixed the same way it prescribes: work
% around it in the fragment, do not touch the script, note it here.
\draw[thinflow] (2.33,1.75) -- (2.86,1.75);

% ---------------- 3. summing node, bias (above) and command (below) ----
\node[mathlab,circle,draw=black!60,line width=0.5pt,fill=white,
      inner sep=0pt,minimum size=3.2mm] (Bsum) at (3.02,1.75) {$+$};
% bias branch, from above -- kept small (cell 0.10) and labelled to its
% WEST (not above) so the column's own top can stay well below the lifting
% arc that passes over it (see step 6): with this column, its top sits at
% y=2.39, and the arc/nets are kept at y>=2.58, a 0.19cm clearance, comfortably
% above the plan's >=0.15cm minimum.
\MatBlock{(2.97,2.39)}{4}{1}{0.10}{black}
\node[mathlab,anchor=east] at (2.94,2.19) {$\symBias$};
\draw[thinflow] (3.02,1.99) -- (3.02,1.91);
% command branch, from below: still the point of the figure, so it keeps
% the full-size cell (0.22) and the accent colour. Column is 4x1 (matches
% B's real shape, (L+1)x1). Label placed to the column's WEST (like b's,
% above) rather than stacked below it -- the first attempt stacked $B$'s
% label directly above the "command" minicap at the old (3-row) offsets,
% and re-using those same y's for a taller 4-row column left only a hair
% of clearance, so the two merged into unreadable "comBand" text.
\MatBlock{(2.91,1.25)}{4}{1}{0.22}{ours}
\node[mathlab,anchor=east] at (2.88,0.81) {$\symBcol$};
\ScalarCell{(3.15,1.03)}{ours}{$\symCmd$}
\draw[oursflow] (3.02,1.25) -- (3.02,1.59);
\node[minicap,anchor=south,text=ours] at (3.02,0.15) {command};

% ---------------- 4. predicted state shat_{t+1} -------------------------
\draw[thinflow] (3.18,1.75) -- (3.70,1.75);
\MemColV{(3.70,2.09)}{4}{0.22}
\node[mathlab,anchor=west] at (3.94,2.05) {$\symMemHat$};

% ---------------- 5. readout: selection row C, then the scalar yhat ----
\draw[thinflow] (3.92,1.75) -- (4.45,1.75);
\RowSelector{(4.45,1.86)}{4}{0.22}
\node[mathlab,anchor=north] at (4.89,1.62) {$\symSel$};
\draw[thinflow] (5.33,1.75) -- (5.86,1.75);
\ScalarCell{(5.86,1.86)}{tok1}{$\symYhat$}

% ---------------- 6. optional learned lifting: one faded detour --------
% 3.1 P6: "the same affine step can be written on learned features E_phi(s_t)
% instead of on the window; this paper reports the identity lifting." Drawn
% once, as a complete arc from s_t's own column top back to shat_{t+1}'s own
% column top, through a de-emphasised encoder/decoder pair -- never a bare
% "id" bypass (that macro belonged to the retired xi-space chain and is not
% used here).
%   The nets sit in the two GAPS either side of the main chain's busy
% middle (Enc over the s_t->K arrow, Dec just before shat_{t+1}), not
% directly above K or the bias column -- the first attempt put Enc right
% above K and its label collided with K's own (L+1)x(L+1) brace label,
% because both are centred on K's own x-span; here Enc's centre (x~1.18)
% is well clear of the brace (x 1.45-2.33). Only the thin horizontal
% connecting line, not a net's full height, has to clear the bias column
% (top y=2.39) and K's brace (top y~2.36) in between -- at y=2.70 it clears
% both by >=0.3cm, comfortably past the plan's 0.15cm minimum.
\begin{scope}[opacity=0.45]
  \draw[bypass] (0.81,2.09) to[out=70,in=200] (1.04,2.80);
  \begin{scope}[shift={(1.10,2.80)},scale=0.24,transform shape]
    \EncNet{mbLiftE}{(0,0)}
  \end{scope}
  \draw[bypass] (1.378,2.80) -- (3.314,2.80);
  \begin{scope}[shift={(3.372,2.80)},scale=0.24,transform shape]
    \DecNet{mbLiftD}{(0,0)}
  \end{scope}
  \draw[bypass] (3.65,2.80) to[out=-20,in=100] (3.81,2.09);
\end{scope}
% E_phi appears exactly once in the whole figure: here, under the one
% de-emphasised encoder that draws it (G10). The decoder that follows it is
% deliberately left unlabelled -- 3.1 P6 introduces one new symbol (E_phi),
% not a paired decoder symbol, and the main text never names one.
%   Two earlier revisions tried to clear the K-dimension label
% ($(L+1)\times(L+1)$, centred over K) by moving THIS label up and by
% shrinking/raising the nets -- both still rendered as merged run-on text
% (only visible in the PNG, not from the coordinates: the brace macro's
% decoration bumps its text noticeably higher than the raw brace-line y
% suggests). Fixed instead by getting this label off that shared row
% entirely: it now sits to the WEST of the encoder, vertically centred on
% the net (y=2.80), far above the brace's own row (y<=2.45).
\node[mathlab,anchor=east] at (0.98,2.80) {$\symEnc$};
\node[minicap,anchor=south west] at (4.55,1.05) {optional learned lifting};
